\\documentclass[11pt]{article}
\\usepackage{amsmath, amssymb, amsthm}
\\usepackage{palatino}

\\title{Inverse Exponentiation as the Next Rung in the Number-System Ladder}
\\author{Thomas Bradley}
\\date{}

\\theoremstyle{definition}
\\newtheorem{definition}{Definition}
\\newtheorem{example}{Example}

\\theoremstyle{plain}
\\newtheorem{theorem}{Theorem}
\\newtheorem{proposition}{Proposition}

\\begin{document}
\\maketitle

\\section{From Exponentiation to Super-Roots}

The historical ``ladder of creation'' for number systems can be summarised as follows:

\\[
\\mathbb{N} \\xrightarrow{\\text{subtraction}} \\mathbb{Z}
  \\xrightarrow{\\text{division}} \\mathbb{Q}
  \\xrightarrow{\\text{roots}} \\mathbb{R}
  \\xrightarrow{\\text{imaginary roots}} \\mathbb{C}.
\\]

At each stage, inverting the dominant operation forces the introduction of a larger number system. We now show that the inverse of exponentiation (super-root or ``power log'') requires leaving~$\\mathbb{C}$.

\\section{A Hole in \\texorpdfstring{$\\mathbb{C}$}{C} for the Super-Root}

\\begin{theorem}
The equation $z^{\\,z} = 0$ has no solution in the complex field $\\mathbb{C}$. Consequently, $\\mathbb{C}$ is not closed under the second-order super-root.
\\end{theorem}

\\begin{proof}
For $z \\in \\mathbb{C}\\setminus\\{0\\}$ we have $z^{\\,z} = e^{\\,z \\log z}$, where $\\log$ denotes the principal complex logarithm. The exponential map never attains~$0$, so no non-zero $z$ satisfies $z^{\\,z}=0$.

At the point $z=0$ we study the limit. Compute
\\[
\\lim_{z \\to 0} z^{\\,z} = \\lim_{z \\to 0} e^{\\,z \\log z} = \\exp\\!\\Big(\\lim_{z \\to 0} z \\log z\\Big) = \\exp(0) = 1,
\\]
where the middle limit evaluates to $0$ by L'H\\^opital's rule. Thus no value of $z \\in \\mathbb{C}$ yields $z^{\\,z} = 0$.
\\end{proof}

\\section{A Quaternion Witness}

\\begin{proposition}
The inverse problem $i^{\\,x} = k$ has no solution $x \\in \\mathbb{C}$ but admits the quaternion solution $x = j$, where $i,j,k$ are the standard quaternion basis elements.
\\end{proposition}

\\begin{proof}
For $z = a + bi \\in \\mathbb{C}$ we have
\\[
i^{\\,z} = e^{\\,z \\log i}
       = \\exp\\!\\left((a+bi)\\,\\frac{i\\pi}{2}\\right)
       = e^{-b\\pi/2}\\big(\\cos(a\\pi/2) + i \\sin(a\\pi/2)\\big),
\\]
which always lies in $\\mathbb{C}$. Hence $i^{\\,z}$ can never equal the purely quaternionic unit~$k$.

In the quaternion algebra, extend exponentiation via
\\[
a^{\\,q} := \\exp\\!\\big((\\log a)\\,q\\big).
\\]
Taking the principal branch $\\log i = i\\pi/2$ and quaternion $j$, we obtain
\\[
i^{\\,j} = \\exp\\!\\left(\\left(\\frac{\\pi}{2} i\\right) j\\right)
       = \\exp\\!\\left(\\frac{\\pi}{2} (i j)\\right)
       = \\exp\\!\\left(\\frac{\\pi}{2} k\\right)
       = k.
\\]
Thus $x=j$ solves $i^{\\,x} = k$, and $x$ lies outside $\\mathbb{C}$.
\\end{proof}

\\section{Conclusion}

The theorem shows a super-root ``hole'' in $\\mathbb{C}$, while the proposition provides an explicit quaternionic solution to a super-root equation. Together they extend the number-system ladder: inverting exponentiation naturally forces the move from $\\mathbb{C}$ into the quaternions~$\\mathbb{H}$.
\\end{document}

